% © 2026 Ing. David Jaroš — CC BY-NC-ND 4.0
%
% This work is licensed under a Creative Commons Attribution-NonCommercial-NoDerivatives
% 4.0 International License (CC BY-NC-ND 4.0).
%
% R ≈ 1.114 — Geometric and Renormalization Identification (B1 + B3)
% Track: CORE — α derivation — v56 baseline
% Date: 2026-03-07

\ifdefined\INCLUDEMODE
  % Being included in another document — skip preamble
\else
  \documentclass[11pt,a4paper]{article}
  \usepackage{amsmath,amssymb,amsthm}
  \usepackage{hyperref}
  \usepackage{booktabs}

  \newtheorem{theorem}{Theorem}
  \newtheorem{proposition}{Proposition}
  \newtheorem{conjecture}{Conjecture}
  \theoremstyle{remark}
  \newtheorem{remark}{Remark}

  \title{$R \approx 1.114$: Geometric Identification\\
    \large Approaches B1 (Volume Ratios) and B3 (Two-Loop on $\psi$-Circle)}
  \author{Ing.\ David Jaroš}
  \date{2026-03-07}

  \begin{document}
  \maketitle
\fi

\section{Problem Statement}
\label{sec:r_problem}

The fine-structure constant derivation in UBT requires
\begin{equation}
  B = B_{\mathrm{base}} \times R
    = N_{\mathrm{eff}}^{3/2} \times R
    \approx 41.57 \times 1.114
    \approx 46.3.
\end{equation}
The correction factor $R \approx 1.114$ is currently phenomenological
(\textbf{OPEN PROBLEM B}).  This document explores two new approaches:

\begin{description}
  \item[B1] $R$ as a ratio of compactification volumes.
  \item[B3] $R$ as a two-loop renormalisation correction on the compact $\psi$-circle.
\end{description}

All calculations avoid using $n^* = 137$ or $\alpha = 1/137$ as inputs
(non-circularity requirement).

% ─────────────────────────────────────────────────────────────────────────────
\section{Approach B1: Volume Ratio Hypothesis}
\label{sec:B1_volume}

\textbf{Status: [DEAD END]} — no volume ratio of natural compactification
spaces reproduces $R \approx 1.114$ without fine-tuning.

\subsection{Tested Combinations}

Let $\mathrm{Vol}(M)$ denote the standard round volume of a manifold $M$.

\begin{center}
\begin{tabular}{llll}
  \toprule
  Formula & Value & Error vs.\ $1.114$ & Status \\
  \midrule
  $\mathrm{Vol}(S^3)/\mathrm{Vol}(S^2\times S^1)$
    & $2\pi^2 / (4\pi \cdot 2\pi) = 1/4$
    & $-78\%$ & FAIL \\
  $\mathrm{Vol}(S^3)/\mathrm{Vol}(\mathbb{RP}^3)$
    & $2\pi^2 / \pi^2 = 2$
    & $+79\%$ & FAIL \\
  $\mathrm{Vol}(\mathbb{RP}^3)/\mathrm{Vol}(S^2)$
    & $\pi^2 / (4\pi) = \pi/4$
    & $-29\%$ & FAIL \\
  Modular volume $\mu(\Gamma_0(7)) = 8$
    & (integer, not dimensionless ratio to $\sim 1$)
    & --- & FAIL \\
  $\mathrm{Vol}(S^3) / (2\pi^2 \cdot \pi)$
    & $2\pi^2 / (2\pi^3) = 1/\pi \approx 0.318$
    & $-71\%$ & FAIL \\
  $\bigl(\mathrm{Vol}(S^1_R) / \mathrm{Vol}(S^1_\psi)\bigr)^{1/N_{\mathrm{eff}}}$
    & (depends on unknown $R$-ratio) & --- & OPEN \\
  \bottomrule
\end{tabular}
\end{center}

None of the natural compactification volume ratios yields a value in the
range $[1.09, 1.14]$.  This approach is therefore classified as
\textbf{[DEAD END]}: the geometric origin of $R$ is not a simple ratio
of sphere volumes.

\begin{remark}
  The modular volume $\mu(\Gamma_0(7)) = 8 = \dim_\mathbb{R}(\mathbb{C}\otimes\mathbb{H})$
  coincides with the real dimension of the biquaternion algebra, but it is an
  integer — not a ratio close to $1.114$.
\end{remark}

% ─────────────────────────────────────────────────────────────────────────────
\section{Approach B3: Two-Loop Renormalisation on the Compact $\psi$-Circle}
\label{sec:B3_twoloop}

\textbf{Status: [OPEN PROBLEM]} — the qualitative structure is plausible and
the required coefficient can be expressed in terms of $N_{\mathrm{eff}}$, but
no combination $f(N_{\mathrm{eff}})$ with a clear algebraic origin reproduces
$R = 1.114$ from first principles.

\subsection{Setup}

The imaginary-time dimension $\psi$ is compactified on a circle
$S^1_\psi$ of radius $R_\psi = \hbar/(m_e c)$.  For a compact spatial
dimension the Kaluza--Klein spectrum is discrete, modifying the
renormalisation group running relative to non-compact flat space.

The one-loop vacuum polarisation on $\mathbb{R}^{1,3}\times S^1_\psi$ reads
\begin{equation}
  \Pi(q^2) = \Pi_{\mathrm{flat}}(q^2)
    + \delta\Pi_{\mathrm{compact}}(q^2),
\end{equation}
where $\delta\Pi_{\mathrm{compact}}$ encodes the finite-size correction from
the discrete KK tower.

\subsection{Ratio of Compact to Flat Renormalisation}

Define the ratio
\begin{equation}
  R_{\mathrm{compact}} = 1 + \frac{\delta\Pi_{\mathrm{compact}}}{\Pi_{\mathrm{flat}}}
    = 1 + \frac{c}{(m_e R_\psi)^2} + \mathcal{O}\!\left(\frac{1}{(m_e R_\psi)^4}\right).
  \label{eq:R_compact_expansion}
\end{equation}
For $R_\psi = \hbar/(m_e c)$ in natural units we have $m_e R_\psi = 1$, so:
\begin{equation}
  R \approx 1 + c, \qquad c \approx 0.114.
\end{equation}

\subsection{Identifying the Coefficient $c$}

From the phenomenological requirement $R = 1.114$:
\begin{equation}
  c = R - 1 = 0.114.
\end{equation}
In a QED-like theory the leading compact-circle correction is
$c \sim \alpha_0 \times f(N_{\mathrm{eff}})$, where
$\alpha_0 = 1/(4\pi B_{\mathrm{base}}) \cdot \text{(normalisation)}$.
Using $\alpha_0 \approx 1/137$ as the leading-order approximation:
\begin{equation}
  f(N_{\mathrm{eff}}) = \frac{c}{\alpha_0} = 0.114 \times 137 \approx 15.6.
\end{equation}

\subsection{Numerical Search for $f(N_{\mathrm{eff}}) \approx 15.6$}

The following combinations of $N_{\mathrm{eff}} = 12$ are evaluated
(see also \texttt{tools/explore\_r\_factor.py}):

\begin{center}
\begin{tabular}{llll}
  \toprule
  Formula for $R$ & Value & Error vs.\ $1.114$ & Status \\
  \midrule
  $1 + \alpha(N_{\mathrm{eff}} + \pi + 1/4)$
    & $1.11235$ & $0.15\%$ & \textbf{[NUMERICAL OBSERVATION]} \\
  $1 + \alpha\,N_{\mathrm{eff}}^{3/2}/e$
    & $1.11162$ & $0.21\%$ & \textbf{[NUMERICAL OBSERVATION]} \\
  $1 + \alpha(N_{\mathrm{eff}} + \pi)$
    & $1.11052$ & $0.31\%$ & \textbf{[NUMERICAL OBSERVATION]} \\
  $1 + \alpha\,N_{\mathrm{eff}} \times (5/4)$
    & $1.10949$ & $0.41\%$ & \textbf{[NUMERICAL OBSERVATION]} \\
  $1 + 1/(3\pi)$
    & $1.10610$ & $0.71\%$ & [CLOSE] \\
  $2^{1/6}$ (sixth root of 2)
    & $1.12246$ & $0.76\%$ & [CLOSE] \\
  \bottomrule
\end{tabular}
\end{center}

The best algebraic candidate from this search is
\begin{equation}
  R \approx 1 + \alpha\bigl(N_{\mathrm{eff}} + \pi + \tfrac{1}{4}\bigr)
    \approx 1.11235,
  \qquad \text{error} = 0.15\%.
  \label{eq:R_best_candidate}
\end{equation}
This is a \textbf{[NUMERICAL OBSERVATION]}.  The structure
$f = N_{\mathrm{eff}} + \pi + 1/4$ can be read as: $N_{\mathrm{eff}} = 12$
discrete bosonic KK modes, $\pi$ from the angular integration over the
half-period of the $\psi$-circle, and $1/4$ a potential two-loop spin
factor.  However, none of these assignments has been derived from
$\mathcal{S}[\Theta]$.

\begin{remark}
  The second-best candidate $1 + \alpha\,N_{\mathrm{eff}}^{3/2}/e \approx 1.11162$
  at $0.21\%$ error is intriguing because $N_{\mathrm{eff}}^{3/2} = B_{\mathrm{base}}$
  connects the $R$ correction directly to $B_{\mathrm{base}}$ via
  $R \approx 1 + \alpha B_{\mathrm{base}}/e$.
  This would imply a self-referential renormalisation structure, but
  requires independent derivation.
\end{remark}

\subsection{Conclusion for B3}

\begin{itemize}
  \item The two-loop compact-circle mechanism provides a \emph{plausible
    qualitative origin} for $R > 1$: finite-$R_\psi$ corrections shift
    the vacuum polarisation upward relative to the flat-space limit.
  \item The best numerical candidate is $R \approx 1 + \alpha(N_{\mathrm{eff}} + \pi + 1/4)$
    at $0.15\%$ error (\textbf{NUMERICAL OBSERVATION}).
  \item No combination has been derived from $\mathcal{S}[\Theta]$ from
    first principles; this approach remains \textbf{[OPEN PROBLEM B]}.
\end{itemize}

% ─────────────────────────────────────────────────────────────────────────────
\section{Numerical Observation from \texttt{explore\_r\_factor.py}}
\label{sec:r_numerical}

Running the companion script (Approach B2, \texttt{tools/explore\_r\_factor.py})
gives the closest algebraic candidate:
\begin{equation}
  R \approx 2^{1/6} = \sqrt[6]{2} \approx 1.1225, \qquad \text{error} = 0.76\%.
  \label{eq:R_2_sixth}
\end{equation}
This is a \textbf{[NUMERICAL OBSERVATION]}: the sixth root of 2 is within
$0.76\%$ of $1.114$.  The number 6 could arise as
$\dim_\mathbb{R}(\mathrm{Im}(\mathbb{C}\otimes\mathbb{H})) - 1 = 7 - 1 = 6$,
or as $N_{\mathrm{phases}} \times N_{\mathrm{helicity}} = 3 \times 2 = 6$.
However, no algebraic derivation of $R = 2^{1/6}$ from
$\mathcal{S}[\Theta]$ exists at this time.

% ─────────────────────────────────────────────────────────────────────────────
\section{Summary}
\label{sec:r_summary}

\begin{center}
\begin{tabular}{lll}
  \toprule
  Approach & Result & Status \\
  \midrule
  B1: Volume ratios               & No match in $[1.09, 1.14]$  & [DEAD END] \\
  B2: Algebraic candidates        & $2^{1/6} \approx 1.1225$, err $0.76\%$ & [NUMERICAL OBSERVATION] \\
  B3: Two-loop on $\psi$-circle   & Correct qualitative structure; $c$ not derived & [OPEN PROBLEM B] \\
  \bottomrule
\end{tabular}
\end{center}

\textbf{Open Problem B remains open.}  The most promising direction is B3
(two-loop compact correction) combined with an algebraic computation of
$c$ from $\mathcal{L}[\Theta]$.

\ifdefined\INCLUDEMODE\else
\end{document}
\fi
